\chapter*{Résumé}
\markboth{Résumé}{}
\begin{singlespace}

Ce rapport présente la conception et la réalisation, durant un stage professionnel d'un mois, d'une application web interne dédiée à la gestion des demandes techniques. Le projet couvre l'enregistrement, la qualification, l'affectation, le traitement, la résolution et la clôture des demandes, avec une historisation des changements importants. L'architecture associe un frontend Angular, une API REST Spring Boot et une base PostgreSQL, selon une organisation modulaire séparant présentation, application, domaine et persistance. La sécurité repose sur l'authentification JWT, le contrôle d'accès fondé sur les rôles et des contrôles métier contextuels. Une assistance IA intégrée via Groq propose une catégorie et une priorité à partir du titre et de la description ; le Responsable technique conserve la décision finale. La validation s'appuie sur des tests automatisés backend et frontend, une intégration locale du parcours IA et une vérification fonctionnelle humaine. Le rapport distingue les résultats mesurés sur le périmètre V1 des limites, notamment le caractère exploratoire du benchmark IA et l'absence de portail client.

\textbf{Mots-clés :} génie logiciel ; API REST ; Angular ; Spring Boot ; intelligence artificielle.

\section*{Abstract}

This report presents the design and implementation of an internal web application for managing technical requests. Developed during a one-month professional internship, the solution addresses a controlled internal scope. The project covers request registration, qualification, assignment, processing, resolution and closure, with the history of significant changes. The architecture combines an Angular frontend, a Spring Boot REST API and a PostgreSQL database, using modular responsibilities for presentation, application, domain and persistence. Security relies on JWT authentication, role-based access control and contextual business checks. An AI assistance feature integrated through Groq proposes a category and a priority from the request title and description; the Technical Manager remains responsible for the final decision. Validation combines automated backend and frontend tests, a local integration of the AI workflow and human functional verification. The report distinguishes measured results within the V1 scope from remaining limitations, including the exploratory nature of the AI benchmark and the absence of a customer portal.

\textbf{Keywords:} software engineering; REST API; Angular; Spring Boot; artificial intelligence.
\end{singlespace}
